% Черновик статьи. Шаблон: https://github.com/sadekovlar/misis-article-latex
% Компилятор: XeLaTeX. Файл misisarticle.cls должен лежать рядом.
% Места, помеченные [ЗАПОЛНИТЬ: ...], требуют ваших данных.
\documentclass{misisarticle}
\usepackage{amsmath,amssymb}

% ГРНТИ/ASJC/OECD — проверьте коды по классификаторам
\grnti{28.23.37, 28.23.19}
\asjc{1702, 1710}
\oecd{1020}

\title{Гибридный метод классификации научно-технических текстов на основе слияния признаков RuBERT и эмбеддингов графа знаний}

\misisauthor{А.~В.~Луценко}{кафедра инженерной кибернетики НИТУ «МИСиС»}{Москва, Россия}{[ЗАПОЛНИТЬ: email]@edu.misis.ru}
\misisauthor{С.~С.~Мишуров}{кафедра инженерной кибернетики НИТУ «МИСиС»}{Москва, Россия}{[ЗАПОЛНИТЬ: email]}

\setabstract{Предложен гибридный метод классификации русскоязычных научно-технических текстов, в котором контекстное представление документа, полученное моделью RuBERT, объединяется на уровне признаков (feature-level fusion) с графовым представлением, вычисленным по заранее обученным эмбеддингам узлов графа знаний. В отличие от подходов типа K-BERT и ERNIE, метод не модифицирует механизм самовнимания трансформера и не требует графовых нейронных сетей при инференсе: извлечение графовых признаков сводится к операции поиска в таблице эмбеддингов со сложностью $O(1)$ на сущность. Формализованы текстовый и графовый потоки, механизм слияния конкатенацией с последующей классификацией многослойным перцептроном, а также протокол экспериментальной оценки по метрикам Macro~F1, Micro~F1, Accuracy и времени инференса. [ЗАПОЛНИТЬ: 1--2 предложения с количественными результатами экспериментов.]}
\setkeywords{классификация текстов, научно-технические тексты, графы знаний, RuBERT, эмбеддинги графа знаний, TransE, слияние признаков}

\begin{document}
\maketitle

\section{Введение}
Архитектура трансформеров стала основой современных методов обработки естественного языка: предобученные модели BERT~\cite{devlin}, SciBERT~\cite{scibert} и их русскоязычные аналоги, в частности RuBERT~\cite{kuratov}, показывают высокое качество в задачах классификации текстов. Вместе с тем научно-технические тексты (НТТ) обладают рядом особенностей: узкоспециализированной терминологией, редкими сущностями, омонимией и сложными концептуальными связями, которые слабо отражаются в статистике совместной встречаемости слов. Трансформер опирается на эти статистические закономерности и не имеет явного доступа к структурированным знаниям предметной области, например к тому, что один метод является разновидностью другого или применяется для определённой задачи.

Естественным решением служит интеграция графов знаний. Однако существующие архитектуры, встраивающие знания непосредственно в механизм внимания, плохо подходят для длинных документов: они увеличивают длину входной последовательности и вычислительные затраты, а ошибки связывания сущностей напрямую искажают представление текста. Альтернативой является разделение обработки на два независимых потока с объединением уже готовых векторных представлений.

\textbf{Цель работы} --- разработка гибридного метода классификации НТТ, объединяющего контекстные эмбеддинги RuBERT с предвычисленными эмбеддингами узлов графа знаний. \textbf{Гипотеза}: слияние указанных признаков на уровне латентных представлений статистически значимо повышает Macro~F1 по сравнению с изолированной моделью RuBERT, при этом время инференса остаётся ниже, чем у методов с глубокой интеграцией графа.

Дальнейшее изложение организовано следующим образом. В разделе~II рассмотрены существующие подходы, в разделе~III описана математическая модель метода, в разделе~IV --- протокол экспериментальной оценки, в разделе~V --- результаты.

\section{Обзор существующих подходов}
\subsection{Специализированные языковые модели}
Применение моделей, обученных на текстах общего назначения, к научным текстам сопряжено со сдвигом предметной области (domain shift). Для его уменьшения была предложена модель SciBERT~\cite{scibert}, предобученная на корпусе из 1,14 млн полных текстов научных статей со специализированным словарём SciVocab, что позволило в среднем на 2--3\% повысить F1 в ряде научных задач по сравнению с BERT. Для русского языка аналогичную роль играет RuBERT~\cite{kuratov}, превосходящий мультиязычные модели на русскоязычных задачах, что обосновывает его выбор в качестве текстового энкодера. Тем не менее любая из этих моделей остаётся статистической и не использует явные онтологические связи между концептами.

\subsection{Интеграция графов знаний в трансформеры}
В K-BERT~\cite{kbert} тройки «субъект--отношение--объект» из графа знаний добавляются к предложению, превращая последовательность в дерево; для ограничения «шума знаний» используются мягкие позиционные эмбеддинги и матрица видимости в слое самовнимания. Метод не требует предобучения с нуля, но увеличивает длину последовательности, усиливает квадратичный рост затрат и чувствителен к ошибкам связывания сущностей, что критично для длинных НТТ.

Модель ERNIE~\cite{ernie} интегрирует знания на этапе предобучения: токены выравниваются с эмбеддингами сущностей, а информация об упоминаниях сущностей вводится в представления токенов. Подход использует сущности как независимые векторы и не учитывает структуру графа явно, а также наследует ограничение длины входа в 512 токенов.

\subsection{Эмбеддинги графов знаний}
Модель TransE~\cite{transe} трактует отношение как трансляцию в вещественном пространстве, $\mathbf{h}+\mathbf{r}\approx\mathbf{t}$, и обучается с помощью ранжирующей функции потерь с зазором. Её известное ограничение --- неадекватное описание отношений типов 1-to-N, N-to-1 и N-to-N. Модель RotatE~\cite{rotate} моделирует отношения как повороты в комплексном пространстве, $\mathbf{t}=\mathbf{h}\circ\mathbf{r}$, $|r_i|=1$, что позволяет выражать симметричность, антисимметричность, инверсию и композицию. Алгоритм node2vec~\cite{node2vec} строит представления узлов по управляемым случайным блужданиям с параметрами $p$ и $q$ и последующим применением Skip-gram, не требуя явных типов отношений.

\subsection{Слияние признаков}
В таксономии мультимодального обучения~\cite{multimodal} выделяют раннее слияние, позднее слияние и слияние на уровне признаков. Последнее позволяет сохранить предобученные пространства признаков каждой модальности и обучить классификатор нелинейным взаимодействиям между ними. Эмпирическое подтверждение дала работа~\cite{ostendorff}: добавление к представлению документа эмбеддингов сущностей Wikidata (TransE, ComplEx) улучшает классификацию по сравнению с BERT без модификации внутренних механизмов трансформера. Наш метод развивает эту идею для русскоязычных НТТ и предметно-ориентированного графа знаний.

\section{Математическая модель метода}
\subsection{Постановка задачи и архитектура}
Требуется построить отображение $f: D\to Y$, где $D$ --- множество документов, $Y=\{y_1,\dots,y_C\}$ --- множество классов (рубрик). Метод состоит из двух потоков и классификатора. Эмбеддинги узлов графа вычисляются заранее (офлайн), поэтому при классификации динамический обход графа не выполняется.

\subsection{Текстовый поток}
Документ токенизируется словарём RuBERT; к последовательности добавляются служебные токены: $X=[x_{\mathrm{CLS}},x_1,\dots,x_n,x_{\mathrm{SEP}}]$, $n+2\le 512$. Модель формирует матрицу скрытых состояний
\begin{equation}
H=\mathrm{RuBERT}(X)\in\mathbb{R}^{(n+2)\times d_{\mathrm{text}}},
\end{equation}
$d_{\mathrm{text}}=768$, а вектор документа берётся из позиции токена [CLS]:
\begin{equation}
V_{\mathrm{text}}=H[0]\in\mathbb{R}^{768}.
\end{equation}
Параллельно NER-модуль выделяет множество научно-технических концептов $E=\{e_1,\dots,e_k\}$.

\subsection{Графовый поток}
Граф знаний --- направленный мультиграф $G=(V_{kg},R)$. Функция вложения $\Phi: V_{kg}\to\mathbb{R}^{d_{kg}}$ обучается заранее на всём графе (TransE, RotatE или node2vec) и далее фиксируется. Каждая сущность $e_i\in E$ связывается с узлом $v_i\in V_{kg}$ (entity linking); несвязанные сущности отбрасываются. Графовое представление документа получается усреднением:
\begin{equation}
V_{\mathrm{graph}}=\frac{1}{|E_{\mathrm{linked}}|}\sum_{e_i\in E_{\mathrm{linked}}}\Phi(v_i),
\end{equation}
а при $E_{\mathrm{linked}}=\varnothing$ принимается $V_{\mathrm{graph}}=\mathbf{0}\in\mathbb{R}^{d_{kg}}$, и классификация опирается только на текстовый поток. Усреднение не зависит от порядка упоминаний, а поиск вектора в таблице имеет сложность $O(1)$ на сущность.

\subsection{Слияние и классификация}
Векторы объединяются конкатенацией:
\begin{equation}
V_{\mathrm{final}}=[V_{\mathrm{text}}\oplus V_{\mathrm{graph}}]\in\mathbb{R}^{d_{\mathrm{text}}+d_{kg}}.
\end{equation}
Классификатор --- перцептрон с одним скрытым слоем:
\begin{equation}
h=\mathrm{ReLU}(W_1V_{\mathrm{final}}+b_1),
\end{equation}
\begin{equation}
\hat{y}=\mathrm{Softmax}\bigl(W_2\,\mathrm{Dropout}(h,p)+b_2\bigr),
\end{equation}
где $W_1\in\mathbb{R}^{d_{\mathrm{hidden}}\times(d_{\mathrm{text}}+d_{kg})}$, $W_2\in\mathbb{R}^{C\times d_{\mathrm{hidden}}}$, $p\in[0{,}3;\,0{,}5]$. Обучение минимизирует перекрёстную энтропию
\begin{equation}
\mathcal{L}=-\sum_{c=1}^{C}y_c\log\hat{y}_c
\end{equation}
оптимизатором AdamW; параметры $W_1,b_1,W_2,b_2$ и, опционально, верхние слои RuBERT дообучаются, а таблица $\Phi$ остаётся замороженной.

\subsection{Вычислительные свойства}
Графовый поток добавляет к стоимости инференса $O(k)$ операций поиска и $O(k\,d_{kg})$ операций усреднения, не зависящих от размера графа и не требующих его обхода. Следует отметить, что стоимость самого трансформера $O(n^2)$ остаётся и входит в оба сравниваемых варианта; выигрыш относится к отсутствию дополнительных затрат на инжекцию троек и динамическую агрегацию соседей. Кроме того, сущности для графового потока могут извлекаться из всего документа, тогда как текстовый поток ограничен окном в 512 токенов.

\section{Протокол экспериментальной оценки}
\subsection{Данные и граф знаний}
[ЗАПОЛНИТЬ: названия датасетов русскоязычных НТТ, число документов и классов, разбиение train/val/test; источник и размер предметно-ориентированного графа знаний (число узлов, отношений); NER-модель и метод entity linking; доля документов, для которых найдена хотя бы одна сущность.]

\subsection{Метрики}
Так как рубрики научных публикаций распределены несбалансированно, основной метрикой выбрана Macro~F1 (чувствительна к редким направлениям); дополнительно применяются Micro~F1 и Accuracy. Эффективность оценивается средним временем инференса на документ (мс/док).

\subsection{Базовые модели}
Сравнение проводится с (1)~TF-IDF с логистической регрессией или SVM; (2)~изолированной моделью RuBERT (только вектор [CLS]); (3)~опционально --- с методами глубокой интеграции графа типа K-BERT и ERNIE. Метод считается успешным, если Macro~F1 превышает значение для изолированного RuBERT при времени инференса не выше, чем у методов глубокой интеграции. Кроме того, сравниваются варианты эмбеддингов графа (TransE, RotatE, node2vec) и значимость различий [ЗАПОЛНИТЬ: статистический тест, число запусков с разными seed].

\section{Результаты и обсуждение}
Результаты экспериментов приведены в табл.~\ref{tab:res}.

\begin{table}[t]
\caption{Сравнение методов на тестовой выборке}\label{tab:res}
\centering
\small
\begin{tabular}{lcccc}
\hline
Метод & Macro F1 & Micro F1 & Acc. & мс/док \\
\hline
TF-IDF + LR & -- & -- & -- & -- \\
RuBERT (text-only) & -- & -- & -- & -- \\
K-BERT / ERNIE & -- & -- & -- & -- \\
Предложенный (TransE) & -- & -- & -- & -- \\
Предложенный (RotatE) & -- & -- & -- & -- \\
Предложенный (node2vec) & -- & -- & -- & -- \\
\hline
\end{tabular}
\end{table}

[ЗАПОЛНИТЬ: обсуждение результатов: вклад графа знаний относительно RuBERT, сравнение алгоритмов эмбеддингов, анализ ошибок для редких классов, влияние качества entity linking, время инференса.]

\section{Заключение}
Предложен и формализован гибридный метод классификации научно-технических текстов, объединяющий контекстные эмбеддинги RuBERT и предвычисленные эмбеддинги узлов графа знаний на уровне признаков. Метод не модифицирует внутренние механизмы трансформера и не требует графовых нейронных сетей при инференсе. [ЗАПОЛНИТЬ: основные количественные выводы.] Дальнейшие исследования предполагают сравнение TransE, RotatE и node2vec, взвешенное (например, на основе внимания) агрегирование эмбеддингов сущностей вместо простого усреднения и анализ устойчивости к ошибкам связывания сущностей.

\begin{thebibliography}{99}
\bibitem{devlin} Devlin J., Chang M.-W., Lee K., Toutanova K. BERT: Pre-training of Deep Bidirectional Transformers for Language Understanding // Proc. NAACL-HLT. 2019. P.~4171--4186.
\bibitem{scibert} Beltagy I., Lo K., Cohan A. SciBERT: A Pretrained Language Model for Scientific Text // Proc. EMNLP-IJCNLP. Hong Kong: ACL, 2019. P.~3615--3620. DOI: 10.18653/v1/D19-1371.
\bibitem{kuratov} Куратов Ю.~М., Архипов М.~Я. Адаптация языковой модели BERT для русского языка и оценка ее качества // Компьютерная лингвистика и интеллектуальные технологии: по материалам конф. «Диалог». 2019. Вып.~18 (25). С.~546--556.
\bibitem{kbert} Liu W., Zhou P., Zhao Z. et al. K-BERT: Enabling Language Representation with Knowledge Graph // Proc. AAAI Conf. on Artificial Intelligence. 2020. Vol.~34, No.~3. P.~2901--2908. DOI: 10.1609/aaai.v34i03.5681.
\bibitem{ernie} Zhang Z., Han X., Liu Z., Jiang X., Sun M., Liu Q. ERNIE: Enhanced Language Representation with Informative Entities // Proc. ACL. 2019. P.~1441--1451.
\bibitem{transe} Bordes A., Usunier N., Garcia-Duran A. et al. Translating Embeddings for Modeling Multi-relational Data // Advances in Neural Information Processing Systems. 2013. Vol.~26. P.~2787--2795.
\bibitem{rotate} Sun Z., Deng Z.-H., Nie J.-Y., Tang J. RotatE: Knowledge Graph Embedding by Relational Rotation in Complex Space // Proc. ICLR. 2019. URL: https://arxiv.org/abs/1902.10197.
\bibitem{node2vec} Grover A., Leskovec J. node2vec: Scalable Feature Learning for Networks // Proc. 22nd ACM SIGKDD. New York: ACM, 2016. P.~855--864. DOI: 10.1145/2939672.2939754.
\bibitem{multimodal} Baltrušaitis T., Ahuja C., Morency L.-P. Multimodal Machine Learning: A Survey and Taxonomy // IEEE Trans. Pattern Anal. Mach. Intell. 2019. Vol.~41, No.~2. P.~423--443. DOI: 10.1109/TPAMI.2018.2798607.
\bibitem{ostendorff} Ostendorff M., Bourgonje P., Berger M., Moreno-Schneider J., Rehm G., Gipp B. Enriching BERT with Knowledge Graph Embeddings for Document Classification // arXiv:1909.08402. 2019.
\end{thebibliography}
\end{document}
